\section{Diagramas de especificação}\label{sec:apendice-diagramas}

É fundamental que todo projeto apresente alguma forma de especificação
do que foi desenvolvido. A descrição é opcional. Assim, \textbf{este
apêndice deve conter os diagramas de especificação que não couberam ao
longo do texto}. Os diagramas devem conter legendas numeradas na
sequência do artigo.

Cada apêndice deve iniciar em uma nova página.

\clearpage
\section{XXX}\label{sec:apendice-xxx}

Podem ser inseridos outros apêndices no artigo tais como códigos de
implementação, telas de interface, instrumentos de coleta de dados,
entre outros. \textbf{Apêndices são} \textbf{textos elaborados pelo
autor} a fim de complementar sua argumentação. Os apêndices são
identificados por letras maiúsculas consecutivas, seguidas de um
travessão e pelos respectivos títulos. Deve haver no mínimo uma
referência no texto anterior para cada apêndice. Colocar sempre um
preâmbulo no apêndice. Caso existam tabelas ou ilustrações,
identifique-as através da legenda, seguindo a numeração normal das
legendas do artigo.
